\lesson{1}{Feb 03 2022 Thu (07:18:18)}{Fears, Anxieties, and Phobias, Oh My!}{Unit 3}

\subsubsection{Enhancing our Vocabulary}

One study revealed that if a person learned how to use a set of $20$ prefixes,
$15$ roots, and $10$ suffixes, the meaning of over 100,000 words could be
unlocked!

\begin{table}[htpb]
  \centering
  \caption{Prefixes, roots, and suffixes}
  \label{tab:label}

  \begin{tabular}{|l|r|}
    \hline
    \bf{-logy}: study the field of & Biology, Technology \\
    \bf{-ite}: one connected with  & Unite, Satellite \\
    \bf{-ia}: state of, or act & Hysteria, Anesthesia \\
    \bf{-ic}: having to do with; pertaining to & Epidemic, Metric \\
    \bf{-ism}: belief in; practice of & Patriotism, Optimism \\
    \bf{-ly}: characteristic of & Quickly, Experimentally \\
    \bf{-ful}: full of & Careful, Beautiful \\
    \bf{-ist}: one who believes in or is engaged in & Dermatologist, Scientist \\
    \bf{y}: characterized by & Happy, Filthy \\
    \hline
  \end{tabular}
\end{table}

\subsubsection{Improving Our Writing}


Professional writers use techniques we can identify to create suspense. An
author may choose just the right word or vivid description to draw readers into
their stories. Another important way authors craft their stories is through
their syntax. Word order in a sentence can be just as effective as word choice
or imagery. Carefully crafting their words—sentence by sentence—allows writers
to create stories that make us laugh, cry, clench our fists in anger, or feel
actual fear.

In our own writing, we can give sentences variety and create desired effects by
manipulating syntax.

\paragraph{Repetition}

Repeating key words or phrases adds emphasis.

\begin{example}[Repetition]
  \begin{itemize}
    \item They saw my terror! They heard my heartbeat! They knew I was guilty!
  \end{itemize}
\end{example}

\paragraph{Fragments}

Using grammatically incomplete sentences demands the reader's attention.

\begin{example}[Fragments]
  \begin{itemize}
    \item Terrible night. Awful night. Oh, the terror of spicy tacos!
  \end{itemize}
\end{example}

\paragraph{Punctuation Marks}

Punctuation controls the pace of the sentence. Choose dashes, commas, colons,
semicolons, exclamation points, and question marks to manipulate how quickly or
slowly the reader experiences your writing.

\begin{example}[Punctuation Marks]
  \begin{itemize}
    \item Fear—or rather yielding to it and thereby giving power to fear mongers
      ---
    \item cannot and will not keep us from pursuing human rights. Because we
      pause to
    \item listen, because we stop to plan our response, because we look at your
    \item arguments slowly and with care, do not think we fear to act. To those
      who
    \item doubt our resolve, we have only to say this: we shall not rest until
      equality
    \item finds its way to the weak, the old, the suffering, and the poor. Rest
      your
    \item head; tomorrow we renew the fight. Chin up! Eyes wide! The battle
      over, what did we gain? 
  \end{itemize}
\end{example}

Authors also use a variety of phrases to create vivid and varied sentences.
Phrases often add description to sentences and tell readers how, where, when,
and why something happened. Take notes on the following types of phrases, paying
particular attention to how the different phrases can be used to give our
sentences variety.

\paragraph{Prepositional}

Phrases:

\begin{itemize}
  \item groups of words that add information to the sentence by functioning as
    an adjective or adverb
  \item always begin with a preposition
\end{itemize}

\begin{example}[Prepositional]
  \begin{itemize}
    \item Under the bridge lived an old troll.
    \item Hope shines a light into the darkest corner.
    \item I found a lost letter from my father beneath the stack of old
      newspapers.
  \end{itemize}
\end{example}

\paragraph{Participial }

Phrases:

\begin{itemize}
  \item groups of words that modify or describe a noun
  \item they always function as adjectives and begin with words ending in -ing
    or -ed
\end{itemize}

\begin{example}[Prepositional]
  \begin{itemize}
    \item The bell, clanging over and over throughout the night, brought some
      comfort to the frightened villagers.
    \item The bell, clanging over and over throughout the night, brought some
      comfort to the frightened villagers.
  \end{itemize}
\end{example}

\paragraph{Gerundial}

Phrases:

\begin{itemize}
  \item groups of words that act as a noun
  \item gerunds end in -ing
\end{itemize}

\begin{example}[Gerundial]
  \begin{itemize}
    \item Screaming in a movie is not usually appropriate unless it is a scary
      movie.
    \item I love telling ghost stories to my younger sister.
  \end{itemize}
\end{example}

\paragraph{Infinitive}

Phrases

\begin{itemize}
  \item groups of words that begin with to + a verb
  \item they can function as adjectives, nouns, or adverbs
\end{itemize}

\begin{example}[Infinitive]
  \begin{itemize}
    \item My plan to frighten my mom backfired.
    \item I wanted to play a harmless trick to see her jump and scream.
  \end{itemize}
\end{example}

\paragraph{Absolute}

Phrases:

\begin{itemize}
  \item groups of words consisting of a noun or pronoun and their modifiers,
    and/or
  \item participle ending in -ing or -ed, and any other modifiers
  \item they add information to the entire sentence
\end{itemize}

\begin{example}[Absolute]
  \begin{itemize}
    \item The old house terrified the neighborhood children, the front door
      creaking like an open mouth and the vacant windows staring like empty
      eyes.
    \item Each child running past the house without looking at it, they hurried
      as if the house might reach out and grab them.
  \end{itemize}
\end{example}

\newpage
